\section{A fully finite proof that \texorpdfstring{$S_5$}{S5} is not solvable}
\label{sec:s5-nonsolvable}

A common proof invokes the simplicity of $A_5$.  For the present purpose that
is more than necessary.  We prove directly that $A_5$ is perfect, using one
explicit commutator.  This reduces the entire obstruction to calculations in
a group of $120$ permutations.

We compose permutations right-to-left and use
$[g,h]=ghg^{-1}h^{-1}$.

\subsection{The commutator subgroup of \texorpdfstring{$S_5$}{S5}}

The sign map
\[
 \sgn:S_5\longrightarrow\{\pm1\}
\]
has abelian target.  Every commutator lies in its kernel, so
\[
 [S_5,S_5]\subseteq A_5.
 \tag{S1}\label{eq:S5-derived-subset}
\]
On the other hand, a direct calculation gives
\[
 [(23),(12)]=(123).
 \tag{S2}\label{eq:S5-commutator}
\]
The commutator subgroup is normal.  Conjugating~\eqref{eq:S5-commutator} by
permutations shows that every $3$-cycle belongs to $[S_5,S_5]$.
Every even permutation is a product of $3$-cycles: write it as an even number
of transpositions and group those transpositions in pairs.  If two
transpositions share a letter, their product is a $3$-cycle; if they are
disjoint, for example,
\[
 (12)(34)=(123)(234).
\]
Hence $3$-cycles generate $A_5$, and
\[
 [S_5,S_5]=A_5.
 \tag{S3}\label{eq:S5-derived}
\]

\subsection{The commutator subgroup of \texorpdfstring{$A_5$}{A5}}

Both
\[
 a=(23)(45),
 \qquad
 b=(12)(45)
\]
are even.  A second direct calculation gives
\[
 [a,b]=(123).
 \tag{S4}\label{eq:A5-commutator}
\]
Thus $(123)\in[A_5,A_5]$.

Every $3$-cycle is conjugate to $(123)$ by an element of $A_5$.  To see this
constructively, first choose any permutation carrying the ordered triple
$(1,2,3)$ to the ordered triple defining the desired $3$-cycle.  If that
permutation is odd, compose it on the right with $(45)$.  This changes its
parity but leaves the images of $1,2,3$ unchanged.  Therefore an even
conjugator always exists.

The subgroup $[A_5,A_5]$ is normal in $A_5$, so it contains every $3$-cycle.
Since those cycles generate $A_5$,
\[
 [A_5,A_5]=A_5.
 \tag{S5}\label{eq:A5-perfect}
\]

\begin{theorem}[Nonsolvability of $S_5$]
\label{thm:S5-not-solvable}
The symmetric group $S_5$ is not solvable.
\end{theorem}

\begin{proof}
By~\eqref{eq:S5-derived}, the first derived subgroup is $A_5$.  By
\eqref{eq:A5-perfect}, every later derived subgroup is again $A_5$:
\[
 S_5^{(1)}=A_5,
 \qquad
 S_5^{(r)}=A_5\quad(r\ge1).
\]
Since $A_5$ has $60$ elements, the series never reaches the identity group.
\end{proof}

\subsection{Primitive-recursive verification}

There are only $5!=120$ permutations.  Encode a permutation by a five-entry
list, enumerate all $120$ lists that pass the bijection test, and compute:

\begin{enumerate}
\item the $60$ even permutations;
\item all commutators in $S_5$ and their generated subgroup;
\item all commutators in $A_5$ and their generated subgroup;
\item the two explicit identities~\eqref{eq:S5-commutator} and
  \eqref{eq:A5-commutator}.
\end{enumerate}

The supplied certificate script finds
\[
 |[S_5,S_5]|=60,
 \qquad
 [S_5,S_5]=A_5,
 \qquad
 |[A_5,A_5]|=60,
 \qquad
 [A_5,A_5]=A_5.
\]
These calculations are not needed as oracles: they can be reproduced by a
small primitive-recursive function and proved correct by induction on the
bounded enumeration.  The explicit human proof above makes the generated
subgroup certificates particularly short.
